\chapter{THEORETICAL BACKGROUND}

\section{Overview of Sign Language Recognition}
\subsection{Definition and Characteristics}
Sign Language Recognition (SLR) is the process of interpreting and translating sign language into text or speech using computer vision and machine learning techniques. Unlike spoken languages, sign languages rely on multiple modalities simultaneously. A single sign is constructed using manual features (hand shape, palm orientation, movement, and location) and non-manual features (facial expressions, head movements, and body posture). The simultaneous nature of these features makes SLR a complex spatio-temporal modeling problem.

\subsection{Vietnamese Sign Language (VSL)}
Vietnamese Sign Language (VSL) is the primary sign language used by the deaf community in Vietnam. VSL exhibits significant regional dialects, primarily categorized into Hanoi, Haiphong, and Ho Chi Minh City dialects. This linguistic diversity introduces high intra-class variance in VSL datasets, meaning the same semantic word might be performed entirely differently depending on the region. Consequently, building a platform to collect VSL must natively support dialect tracking and metadata tagging, which poses a unique database modeling challenge.

\section{Landmark-Based Representation and MediaPipe}
\subsection{Human Pose Estimation Concept}
Human Pose Estimation (HPE) is a computer vision task that identifies and tracks the anatomical joints (keypoints) of the human body in images or videos. HPE serves as the foundational preprocessing step for modern gesture recognition systems, abstracting away background noise, lighting conditions, and subject clothing to focus solely on movement kinematics.

\subsection{Google MediaPipe Holistic Architecture}
MediaPipe Holistic, developed by Google, is a state-of-the-art framework for simultaneous perception of human body, face, and hand kinematics \cite{lugaresi2019mediapipe}. It integrates distinct machine learning models: a pose estimation model (BlazePose), a face mesh model, and a hand tracking model. 
MediaPipe Holistic provides 543 3D landmarks (33 pose landmarks, 21 per hand, and 468 facial landmarks) per frame. By leveraging lightweight Blaze neural network architectures, MediaPipe can run in real-time within a web browser using WebAssembly (Wasm). This allows the CTU-SignBridge platform to extract features directly on the client's device, significantly reducing the server's computational burden, a technique increasingly adopted in recent sign language recognition literature \cite{gomez2024sign, sharma2024continuous, wijaya2023dynamic}.

\subsection{Data Storage Optimization (NPZ vs. Raw Video)}
A critical challenge in crowd-sourced video collection is storage and bandwidth consumption. A typical 3-second 1080p video recorded at 30 frames per second (FPS) consumes approximately 2 to 5 MB of disk space. When scaling to a dataset of 100,000 samples, storage requirements exceed 500 GB.
By extracting MediaPipe landmarks in the browser, the data representation is transformed from a dense RGB pixel matrix into a sparse tensor of dimensions $(T \times 543 \times 4)$, representing $T$ frames, 543 landmarks, and 4 spatial coordinates ($x, y, z, visibility$). Compressing this tensor using the NumPy Zipped (\texttt{.npz}) format yields a file size of approximately 50 to 100 KB per sample. This represents a 98\% reduction in bandwidth and storage costs compared to raw video, enabling scalable and cost-effective data collection.

\section{Database Modeling and Taxonomy}
\subsection{Relational Database Principles}
Relational Database Management Systems (RDBMS), such as PostgreSQL, organize data into tables with predefined schemas. Relationships between entities are established using primary and foreign keys, ensuring data integrity through ACID (Atomicity, Consistency, Isolation, Durability) properties. In the context of a sign language taxonomy, data integrity is paramount to prevent orphaned vocabulary terms or inconsistent dialect mappings.

\subsection{Limitations of Hierarchical Models}
Traditionally, taxonomies are modeled using hierarchical structures such as the Adjacency List model (where a row contains a \texttt{parent\_id} reference). While simple to implement, hierarchical models are highly inefficient for complex querying. Retrieving an entire vocabulary tree requires recursive Common Table Expressions (CTEs), which degrade database performance significantly as the depth of the taxonomy grows. Furthermore, sign languages often require polyhierarchical categorization (e.g., a sign belongs to both "Action" and "Emotion" categories), which strict tree models cannot efficiently represent.

\subsection{Dimensional Modeling and Star Schema}
To overcome the limitations of hierarchical models, the Star Schema -- a concept derived from Data Warehousing -- is employed. In a Star Schema, a central \textit{Fact Table} (representing the core entity, such as the Sign Vocabulary) is surrounded by denormalized \textit{Dimension Tables} (such as Category, Region, and Handshape). 
This schema design eliminates the need for deep recursive joins. Instead, it allows for high-performance, single-level `JOIN` operations, making it extremely efficient for multi-faceted filtering (e.g., "Find all signs in the 'Food' category using the 'Southern' dialect").

\section{Multi-Tenancy Architecture}
\subsection{Cloud Computing and SaaS Concepts}
Software as a Service (SaaS) is a software distribution model where applications are hosted by a cloud provider and made available over the internet \cite{schroeter2012enabling}. A fundamental architectural requirement for SaaS is \textbf{Multi-Tenancy}, wherein a single instance of the software and its supporting infrastructure serves multiple distinct customer groups (tenants) \cite{bezemer2010multi, wu2014tossma}. In the CTU-SignBridge platform, a "tenant" corresponds to an independent academic research team or institution.

\subsection{Data Isolation Strategies}
Ensuring strict data isolation between tenants is critical to prevent data leakage and ensure privacy. There are three primary strategies for multi-tenant database isolation:
\begin{enumerate}
    \item \textbf{Database-per-Tenant:} Maximum isolation, highest cost. Each tenant has a separate physical database instance.
    \item \textbf{Schema-per-Tenant:} Moderate isolation and cost. Tenants share a database instance but have separate logical schemas.
    \item \textbf{Shared-Schema (Row-level Isolation):} Lowest cost, highest complexity. All tenants share the same tables, and isolation is enforced programmatically via a \texttt{tenant\_id} (or \texttt{workspace\_id}) column in every row.
\end{enumerate}
CTU-SignBridge adopts the \textbf{Shared-Schema} approach to maximize resource utilization and simplify schema migrations, relying on the application layer and Casbin for strict authorization checks.

\section{Identity and Access Management (IAM)}
\subsection{Authentication vs. Authorization}
Identity and Access Management relies on two distinct concepts:
\begin{itemize}
    \item \textbf{Authentication (AuthN):} The process of verifying \textit{who} a user is (e.g., via username/password or OAuth).
    \item \textbf{Authorization (AuthZ):} The process of determining \textit{what} an authenticated user is permitted to do (e.g., read, write, or delete specific resources).
\end{itemize}

\subsection{Keycloak for Single Sign-On (SSO)}
Keycloak \cite{keycloak2024} is an open-source Identity and Access Management solution developed by Red Hat. It provides robust Authentication mechanisms, including Single Sign-On (SSO), social login, and multi-factor authentication (MFA). By delegating authentication to Keycloak, CTU-SignBridge offloads the complex responsibilities of password hashing, token issuance (JWT), and session management to a secure, enterprise-grade identity provider.

\subsection{Casbin for Policy-Based Access Control (RBAC/ABAC)}
While Keycloak handles Authentication, fine-grained Authorization is managed by Casbin \cite{casbin2024}. Casbin is a powerful and efficient open-source access control library that supports multiple models, including Role-Based Access Control (RBAC) and Attribute-Based Access Control (ABAC). 
In CTU-SignBridge, Casbin enforces policies that dictate tenant boundaries (e.g., user $A$ with role $Annotator$ in $Workspace_1$ cannot access $Workspace_2$). Casbin's PERM (Policy, Effect, Request, Matcher) metamodel allows for highly dynamic and scalable permission evaluations without hardcoding authorization logic into the application's business layer.

\section{Asynchronous Processing and Object Storage}
\subsection{Message Broker Architecture (Redis)}
A Message Broker is an architectural pattern for message validation, transformation, and routing. It facilitates inter-process communication by mediating between a producer (sending a message) and a consumer (receiving the message). \textbf{Redis} \cite{redis2024}, an in-memory data structure store, is utilized in this architecture as an extremely fast, lightweight message broker to queue background tasks.

\subsection{Distributed Task Queues (Celery)}
In a robust web application, long-running processes (such as cloud synchronization, video processing, or bulk email dispatch) should not block the main HTTP thread. \textbf{Celery} \cite{celery2024} is a distributed task queue built on Python that integrates seamlessly with Redis. 
When a user uploads a dataset, the API server responds immediately to the client, while a Celery worker asynchronously processes the file upload in the background. If a network failure occurs during the background task, Celery's built-in retry mechanisms with exponential back-off ensure that the task is automatically re-queued, providing high fault tolerance.

\subsection{S3-Compatible Object Storage (MinIO)}
Unlike relational databases designed for structured text data, binary assets (such as \texttt{.npz} arrays and MP4 videos) require specialized storage solutions known as Object Storage. \textbf{MinIO} \cite{minio2024} is a high-performance, S3-compatible object storage server. By deploying MinIO, CTU-SignBridge mimics the AWS S3 API locally, allowing the platform to store vast quantities of unstructured data reliably. MinIO provides robust features such as data replication, object locking, and seamless migration to public cloud providers (like Amazon S3 or Google Cloud Storage) if future scalability demands it.
